% Part: proof-theory
% Chapter: sequent-calculus
% Section: introduction

\documentclass[../../../include/open-logic-section]{subfiles}

\begin{document}

\olfileid{pt}{seq}{int}

\olsection{مقدمة}

حسابات المتتاليات أنظمة~!!{proof} استحدثها غيرهارد غنتسن في ثلاثينيات
القرن العشرين. ولم يكن المقصود منها في المقام الأول أن تكون طريقة عملية
لـ!!{proving} الأشياء، بل استخدمها غنتسن كي يتمكن من إثبات نتائج معينة
عن البنية الممكنة لـ!!{proof}s وخصائصها. وهذه الأنواع من المسائل هي بالضبط
المسائل التي تعنى بها نظرية البرهان.

تعمل حسابات المتتاليات على \emph{المتتاليات}، كما يوحي الاسم، لا على
!!{formula}s. والمتتالية زوج من تسلسلين أو متعددتي مجموعات من
!!{formula}s. استعمل نظاما غنتسن الأصليان (\Log{LK} و\Log{LJ})
التسلسلات؛ أما منظّرو البرهان فيفضّلون اليوم متعددات المجموعات. ومتعددة
المجموعات تشبه المجموعة، إلا أنها قد تحتوي~!!{element} أكثر من مرة. ومع
ذلك لا يهم ترتيب~!!{element}s فيها، كما في المجموعات. ولذلك يصف التعبيران
$\{!A, !A, !B\}$ و$\{!A, !B, !A\}$ متعددة المجموعات نفسها، غير أن هذه
المتعددة تختلف عن~$\{!A,!B\}$ وعن~$\{!A,!A,!B,!B\}$، لأن هاتين الأخيرتين
تحتويان أعدادًا مختلفة من نسخ~$!A$ و$!B$ عن الأولى.

تُكتب متعددات مجموعات~!!{formula}s في نظرية البرهان عادةً بحروف يونانية
كبيرة. فعندما نكتب~$\Gamma$ أو~$\Delta$ أو~$\Pi$ أو~$\Lambda$ أو~$\Theta$،
نعني أنها متعددات مجموعات من~!!{formula}s في المنطق الذي نعمل فيه. ومن
المتعارف عليه أيضًا أن منظّري البرهان لا يكتبون~$\tuple{\Gamma,\Delta}$،
بل يفصلون المركبتين برمز متتالية؛ ونستعمل نحن~$\Sequent$.

\begin{defn}[متتالية]
الـ\emph{متتالية} تعبير من الصورة
\[
\Gamma \Sequent \Delta
\]
حيث~$\Gamma$ و$\Delta$ متعددتا مجموعات منتهيتان (وقد تكونان فارغتين) من
!!{formula}s. وتسمى~$\Gamma$ \emph{المقدِّمة}، بينما تسمى~$\Delta$
\emph{المؤخِّرة}.
\end{defn}

لتكن~$!G$ اقتران~!!{formula}s الموجودة في~$\Gamma$ (مع اعتبار الاقتران
الفارغ~$\ltrue$)، ولتكن~$!D$ فصل~!!{formula}s الموجودة في~$\Delta$ (مع
اعتبار الفصل الفارغ~$\lfalse$). عندئذ تكون المتتالية~$\Gamma\Sequent\Delta$
صادقة (في بنية ما~$\Struct{M}$) إذا وفقط إذا كانت~$!G\lif !D$ صادقة. وفي
المنطق الكلاسيكي يكافئ هذا القول إن إحدى~!!{formula}s في~$\Gamma$ كاذبة
أو إن إحدى~!!{formula}s في~$\Delta$ صادقة.

إذا كانت~$\Gamma$ متعددة مجموعات من~!!{formula}s، فإننا نكتب~$\Gamma,!A$
(أو~$!A,\Gamma$) لنتيجة إضافة~$!A$ إلى~$\Gamma$. وإذا كانت~$\Delta$ أيضًا
متعددة مجموعات من~!!{formula}s، فإن~$\Gamma,\Delta$ هو اتحاد متعددتي
المجموعات: فإذا كانت~$\Gamma$ تحتوي~$!A$ بـ\emph{تعددية}~$n$ (أي تحتوي~$n$ نسخة
من~$!A$)، وكانت~$\Delta$ تحتوي~$!A$ بتعددية~$m$، فإن~$\Gamma,\Delta$
تحتوي~$!A$ بتعددية~$n+m$. وإذا كانت متعددة المجموعات في أحد جانبي
المتتالية فارغة، فإننا نترك موضعها خاليًا عادةً. فمثلًا،~$\Sequent !A$
هي المتتالية ذات المقدِّمة الفارغة والمؤخِّرة التي تحتوي~$!A$ بتعددية~$1$.

يتكون~!!^a{proof} في حساب المتتاليات من شجرة متتاليات، تكون فيها
المتتاليات العليا (أو الابتدائية) بديهيات النظام المدروس، وإذا وقعت متتالية
تحت متتالية أو اثنتين غيرها وجب أن تنتج منهما إنتاجًا صحيحًا بقاعدة
استدلال في النظام. سننظر أولًا في نظامين مختلفين من حساب المتتاليات للمنطق
الكلاسيكي، هما~$\Log{G1c}$ و$\Log{G3c}$. وترد بديهياتهما وقواعدهما في
\cref{pt:seq:int:tab:G1c,pt:seq:int:tab:G3c}. \oliflabeldef{fol:seq::chap}{
(ويرد نظام غنتسن الأصلي~\Log{LK} في~\olref[fol][seq][]{chap}.)}{}

\subfile{rules-G1c}
\subfile{rules-G3c}

تختلف الأنظمة فيما بينها اختلافًا دقيقًا، وتترتب على هذه الفروق خصائص
برهانية نظرية مختلفة. فبديهيات~\Log{G1c} هي~$!A\Sequent !A$ ولا تسمح
بـ!!{formula}s أخرى. أما بديهيات~\Log{G3c} فمن الصورة~$!C,\Gamma\Sequent
\Delta,!C$، حيث يجوز وجود~!!{formula}s «جانبية» اعتباطية في~$\Gamma$
و$\Delta$، لكن~!!{formula}~$!C$ التي تظهر في الجانبين يجب أن تكون
\emph{ذرية}. ولـ\Log{G1c} نسختان من كل من قاعدتي~$\LeftR{\land}$
و$\RightR{\lor}$، في حين أن لـ\Log{G3c} نسخة واحدة فقط من كل منهما.
ويضم~\Log{G1c} كذلك «قواعد بنيوية» إضافية تسمى «التقلص»
(\LeftR{\Contraction}، \RightR{\Contraction}) و«الإضعاف»
(\LeftR{\Weakening}، \RightR{\Weakening}).

نظاما~\Log{G1c} و\Log{G3c} سليمـان وكاملان للمنطق الكلاسيكي. أي إن كل
متتالية~!!{provable} في هذين النظامين صادقة في كل~!!{structure}، وكل
متتالية صادقة في كل~!!{structure} لها~!!a{proof}. ولذلك يمكن أيضًا الإجابة
عن سؤال \emph{ما إذا كانت} لمتتالية ما~!!a{proof} بطرائق منطقية أخرى (كطرائق
نظرية النماذج). ولأن كلا النظامين~!!{prove} المتتاليات الصادقة في كل
!!{structure} بالضبط، فإنهما يبرهنان المتتاليات نفسها بوجه خاص.

لكن نظرية البرهان لا تعنى فقط بـ\emph{ما إذا كان} شيء ما قابلًا للبرهنة،
بل أيضًا بـ\emph{كيفية} برهنته. وينظر منظّرو البرهان في مسائل من قبيل:
«هل يمكن دومًا تجنب قاعدة معينة؟»، و«هل يمكن إضافة هذه القاعدة إلى النظام
من دون أن يؤدي ذلك إلى قابلية متتاليات جديدة للبرهنة؟»، و«هل يمكن تحويل
!!{proof}s في نظام إلى~!!{proof}s في نظام آخر؟». وإذا كان الجواب عن مسألة
من هذا النوع «نعم»، فكثيرًا ما تُثبت هذه الحقيقة بالاستقراء على تعقيد
!!{proof}s، أو على بنية~!!{proof}s على نحو مكافئ. ولا ينتج من ذلك برهان الحقيقة
فحسب (مثلًا: إذا برهن~\Log{G1c} متتاليةً فإن~\Log{G3c} يبرهن المتتالية
نفسها)، بل ينتج أيضًا إجراء (مثل تحويل~!!{proof}s في~\Log{G1c} إلى
!!{proof}s في~\Log{G3c}).

ومن النتائج المركزية في نظرية البرهان~\emph{مبرهنة حذف القطع}. فهي تبين
أننا نستطيع إضافة قاعدة القطع
\[
\Axiom$\Gamma \fCenter \Delta, !A$
\Axiom$!A, \Pi \fCenter \Lambda$
\RightLabel{\Cut}
\BinaryInf$\Gamma, \Pi \fCenter \Delta, \Lambda$
\DisplayProof
\]
إلى نظامي المتتاليات من دون تغيير المتتاليات~!!{provable}. وفوق ذلك، يقدم
البرهان طريقة تحوّل أي~!!{proof} يستعمل قواعد~\Log{G1c} (أو~\Log{G3c})
مع~\Cut{} إلى برهان لا يستعمل إلا قواعد~\Log{G1c} (أو~\Log{G3c}). وبعبارة
أخرى، «يحذف» هذا الإجراء قاعدة~\Cut{} من~!!{proof}s التي تستعملها؛ ومن هنا
جاء الاسم.

\end{document}
